% 第4章 详细设计与实现
\chapter{TIDE-Swarm详细设计与实现}
在第三章总体架构的基础上，本章介绍TIDE-Anchor上下文连续性与记忆保持、TIDE-Mesh动态异构协作与低熵通信、TIDE-Placement端边云推理调度三项核心机制，说明其状态模型、运行过程与核心算法，并分析计算与通信开销。三项机制通过任务状态和执行反馈衔接，共同支撑长程任务的持续推进与异常恢复。

面向赛题提出的上下文连续性、动态异构拓扑、低熵通信、神经符号协同推理与端边云调度要求，本文以\textbf{语义生成、显式校验、受控执行和证据反馈}组织长程协作。神经侧通过大模型理解自然语言意图、生成任务计划与阶段结果，并在配置质量评审模型时提供语义判断；符号侧将任务目标、依赖、资源需求和验收条件组织为结构化状态，以计划校验、资源硬过滤、消息字段保护和质量规则约束执行。两侧通过计划中间表示、反例、任务账本与结构化消息交换信息，使推理建议转化为可检查、可修订的执行依据。

三项机制在这一闭环中分别回答“依据什么继续推理”“由哪些成员协作”“在哪类资源上执行”。TIDE-Anchor提供稳定目标与可追溯证据，TIDE-Mesh依据当前任务、能力需求和失败反馈生成稀疏的活跃协作关系，TIDE-Placement依据实时性、敏感度和资源状态安排任务阶段。\textbf{神经符号协同贯穿三项机制的决策接口}：语义判断提出候选，显式规则限定可行范围，运行结果再更新计划、成员选择和资源分配。第\ref{sec:mechanism_coordination}节进一步说明计划校验、动态拓扑与分级恢复如何衔接，并给出可观测的评价口径。


% 重构稿：版本化共享状态、四级记忆与四类长程失效治理。
\section{TIDE-Anchor超长程上下文连续性与记忆保持}
\label{sec:tide_anchor}

\subsection{问题定义与总体设计}

长程任务中的上下文连续性，是在任务跨越多轮工具调用、角色交接和执行中断时，持续提供正确目标、当前状态和可核验证据的能力。数小时乃至数天是本机制面向的设计场景，实际持续运行能力仍需按实验时长验证。上下文窗口扩展只能增加信息容量，不能自动解决信息选择、错误隔离和恢复一致性问题。

本文围绕四类相互关联的失效组织TIDE-Anchor。第一，上下文溢出与退化：工具结果持续累积，历史截断可能丢失关键参数和执行结果。第二，指令稀释与注意力漂移：约束虽仍在输入中，却未被当前决策有效利用；长上下文中的位置敏感性已有相关研究讨论\cite{tide_lost_middle}。第三，误差累积与级联传播：未核验的阶段输出被下游当作事实，局部故障转化为后续决策的错误前提。第四，静默置信度失效与状态丢失：服务调用成功却没有任务进展，或中断后无法一致恢复信息与执行位置。

\begin{table}[htbp]
\centering
\caption{四类长程失效与TIDE-Anchor的对应机制}
\label{tab:anchor_failure_mapping}
\begin{tabular}{p{0.21\textwidth}p{0.40\textwidth}p{0.28\textwidth}}
\toprule
失效模式 & 设计机制 & 直接检查对象 \\
\midrule
上下文溢出与退化 & 状态投影、原文归档、四级记忆检索与装配后预算检查 & 输入用量、证据覆盖与归档可访问性 \\
指令稀释与漂移 & 共享状态中的目标和约束每步重锚定，动作前后进行契约检查 & 约束满足、目标关联与待办进展 \\
误差累积与级联 & 候选结果隔离、证据验证、依赖版本追踪与局部重执行 & 已验证状态、来源与下游依赖 \\
静默失效与状态丢失 & 进展检测、统一状态检查点与恢复前一致性校验 & 实际产物变化、执行前沿及副作用记录 \\
\bottomrule
\end{tabular}
\end{table}

TIDE-Anchor以\textbf{版本化共享状态和四级记忆}构成节点信息供给的两个来源。共享状态维护当前任务的权威视图；记忆保存已完成工作提炼出的、具有适用条件的可复用信息。二者经上下文装配器进入节点，原始工具输出和文件保留在归档中。账本与检查点统一归入共享状态管理组件：账本是当前状态的可读视图，检查点是这一状态连同执行位置的可恢复版本。

\begin{figure}[htbp]
\centering
\fbox{\parbox{0.92\textwidth}{\centering
版本化共享状态：目标、约束、进度、已验证事实、执行前沿\\[3pt]
$\downarrow$ 状态投影\qquad\qquad\qquad $\downarrow$ 检索条件\\[3pt]
当前输入 $+$ 任务锚点 $+$ 四级记忆检索证据\\[3pt]
$\downarrow$ 来源核验与预算检查 $\downarrow$\\[3pt]
节点执行 $\longrightarrow$ 结果验证 $\longrightarrow$ 状态提交与检查点\\[3pt]
\hfill 已完成子任务 $\longrightarrow$ 经验总结 $\longrightarrow$ 记忆准入}}
\caption{共享状态与四级记忆协同的上下文供给架构}
\label{fig:anchor_memory}
\label{fig:anchor_context}
\end{figure}

\subsection{共享状态与四级记忆的职责边界}

\paragraph{版本化共享状态}
设第$t$个已提交版本为
\begin{equation}
S_t=\langle g,\mathcal{C}_t,\mathcal{P}_t,\mathcal{F}_t,
\mathcal{A}_t,\mathcal{D}_t,\mathcal{X}_t,v_t\rangle,
\end{equation}
其中，$g$为原始目标，$\mathcal{C}_t$为当前有效约束，$\mathcal{P}_t$为待办及完成状态，$\mathcal{F}_t$为带验证标记的事实，$\mathcal{A}_t$为产物与归档引用，$\mathcal{D}_t$为证据和产物依赖，$\mathcal{X}_t$为执行前沿、预算及尝试记录，$v_t$为版本。节点产物先进入候选区，通过相应验收后再更新权威字段。未通过验证的推测可以保留，但必须携带未验证或争议状态。

共享不意味着把整个状态复制给每个模型。节点只获得其权限和职责允许读取的投影$P_n(S_t)$；原始目标、有效硬约束及本步验收条件构成任务锚点。并行节点基于明确版本读取状态，写回时检查基础版本和字段归并规则，避免后提交的旧结果覆盖更新后的约束。无法自动消解的写入冲突进入协调流程。

\paragraph{四级记忆的产生与使用}
记忆来源于\textbf{已结束的节点操作或已完成子任务的总结}，其作用是复用经验和证据，而不是复制当前任务状态。整个任务尚未结束时，已验收子任务也可以形成TASK记忆；跨任务的PROJECT或GLOBAL记忆则需要更严格的归纳和准入。运行中的草稿、即时工具返回和未经检查的模型输出优先留在工作缓冲或候选结果区，不能直接作为稳定知识传播。

\begin{table}[htbp]
\centering
\caption{四级记忆的产生条件与复用范围}
\label{tab:anchor_memory_scopes}
\begin{tabular}{p{0.15\textwidth}p{0.44\textwidth}p{0.30\textwidth}}
\toprule
作用域 & 形成条件与典型内容 & 读取边界 \\
\midrule
WORKING & 已结束局部操作的短摘要；尚待子任务验收时标为候选 & 当前节点或局部执行链，短生命周期 \\
TASK & 已验收子任务的结论、参数解释和失败经验 & 当前运行，绑定任务版本 \\
PROJECT & 完成任务后经核验的项目约定、处理方法和历史结论 & 同项目跨运行，受权限和有效期约束 \\
GLOBAL & 从完成任务中抽象、去除私有信息并通过准入的通用经验 & 经授权跨项目使用，不携带项目秘密 \\
\bottomrule
\end{tabular}
\end{table}

四级表示使用范围，不是同一条消息逐级自动升级的流水线。记忆应记录适用条件、结论、证据引用、有效时间、来源任务和验证状态；失败经验也可保留，但应表达“何种条件下失败”，不能把失败结果写成正确事实。HOT、WARM、STRUCTURED、SEMANTIC和COLD属于存储介质，与上述作用域独立配置。

\begin{table}[htbp]
\centering
\caption{共享状态、检查点、记忆与归档的关系}
\label{tab:anchor_state_objects}
\begin{tabular}{p{0.20\textwidth}p{0.34\textwidth}p{0.35\textwidth}}
\toprule
对象 & 主要问题 & 更新与恢复语义 \\
\midrule
共享状态当前视图 & 现在要完成什么，哪些结果可信 & 验收后按版本提交，账本呈现该视图 \\
共享状态检查点 & 中断后从哪个一致版本继续 & 保存状态、执行前沿和必要恢复引用 \\
四级记忆 & 哪些已完成工作的信息可在此复用 & 总结后准入，按范围检索，可失效或撤销 \\
原文归档 & 如何核验摘要和阶段结论 & 保存工具结果及产物原文，以引用和摘要校验关联 \\
\bottomrule
\end{tabular}
\end{table}

\subsection{面向上下文溢出与退化的按需装配}

每个节点的上下文由当前输入、共享状态投影和记忆检索结果共同构造：
\begin{equation}
C_{n,t}=O_n\oplus A_n(S_t)\oplus I_{n,t}
\oplus P_n^{\mathrm{detail}}(S_t)\oplus E_{n,t},
\label{eq:anchor_context}
\end{equation}
其中，$O_n$为角色和输出契约，$A_n(S_t)$为任务锚点，$I_{n,t}$包含本步输入及第\ref{sec:tide_mesh}节通信模块交付的定向消息，$P_n^{\mathrm{detail}}$为其余必要状态，$E_{n,t}$为从四级记忆唤醒的证据。状态投影和检索结果按来源标识去重；通信信箱不再形成独立的全量历史副本。

检索先按任务、项目、权限、时间和验证状态限定候选，再结合语义与词项匹配排序。WORKING和TASK优先补充本次运行的局部依据，PROJECT和GLOBAL提供条件适配的历史经验；跨层候选可统一重排。命中摘要后，仅在核对参数、变更或来源时读取原文。Silent、Standard、Deep和Recovery是检索深度配置，分别面向局部操作、常规执行、复核与恢复，不改变权限边界。

\begin{equation}
B_{\mathrm{in}}=\max(0,B_{\mathrm{max}}-B_{\mathrm{out}}-B_{\mathrm{margin}}),
\qquad \operatorname{Tok}(C_{n,t})\leq B_{\mathrm{in}},
\label{eq:anchor_budget}
\end{equation}
其中$B_{\mathrm{margin}}$为工具格式和计数误差预留空间。预算检查针对最终装配结果，而不只检查原始共享状态。超限时依次减少低相关记忆、缩短可回溯证据正文、移出冗余状态描述；目标、硬约束、必要证据引用和输出契约优先保留。若必要信息本身已超限，则暂停或拆分任务，不以静默截断破坏执行条件。

长工具结果在写回时形成摘要与原文引用，完整载荷留在归档。该设计把模型可见输入的增长限制在当前决策所需范围内，但摘要可能遗漏细节，检索也可能漏召回，因此仍需通过证据充分性及关键参数保留率检查实际效果。

\subsection{面向指令稀释与注意力漂移的任务重锚定}

任务重锚定在每步重新构造目标与约束，使它们不依赖历史检索是否命中。锚点包含原始目标、当前有效约束、本步目标、已确认事实、不可假定事项和验收条件。上下文模板把角色与约束置于明确位置，并在临近当前动作处给出简短验收清单；这种输入组织用于减少约束被过程信息淹没的机会，不构成对模型注意力的理论保证。

状态约束与历史经验冲突时，先检查二者的有效版本与来源。当前用户授权和有效任务约束决定执行边界，记忆只能提出参考或冲突线索。系统不能因历史任务曾允许某种行为，就改变本次任务的权限、资源范围或验收条件。

模型执行前，检查候选动作的工具、参数和资源是否满足显式约束；执行后，再将结果与当前验收条件核对。语义目标关联可由评审模型辅助判断，可编码的权限、数值和输出结构采用规则检查。发现偏离时生成定位明确的修订指令；连续偏离或修订预算耗尽时暂停并记录原因。该路径将重锚定与可观测检查连接起来，避免仅依赖反复复述目标。

\subsection{面向误差累积与级联传播的证据隔离与局部修复}

结果隔离以“已返回”和“已确认”分开处理。HTTP成功、工具正常退出或模型自报高置信度只能说明一次调用完成，不能自动证明业务结果正确。节点输出先标为候选，按任务性质检查非空产物、来源、参数一致性、测试或其他验收证据；通过后提交状态，再提炼相应记忆。工具异常与低质量结果保留为失败记录，不填入成功事实字段。

依赖追踪记录节点实际使用的事实、记忆和产物版本。设结论$y$的输入依赖为
\begin{equation}
\operatorname{Dep}(y)=\{(x_i,v_i)\}_{i=1}^{m}.
\end{equation}
当某版本被证实错误或过期时，先隔离该记录并展开原文复核，再将直接及间接依赖它的结果标为待重验。对已失效依赖的后继执行设置阻断，避免错误继续进入第\ref{sec:tide_mesh}节的消息传输与团队交接。

\paragraph{A到B再到C的修复示例}
设A提取数据单位，B据此计算指标，C生成报告。B引用历史记忆$m^{(2)}$，将“元”误读为“万元”。若C在生成报告前检查到指标与原始表格不一致，先暂停依赖该指标的报告步骤，读取证据确认是记忆错误、A的提取错误，还是B的计算错误。若确认只有$m^{(2)}$错误，则撤销该版本，保留A中独立且正确的提取结果，使用核验后的单位重新执行B，再执行C；若错误来自A，则A及受影响的B、C均需重做。单纯检索得分低或模型不确定时，应先补证，不能直接认定历史记忆错误。

若修复期间用户增加了新约束，旧检查点只作为重建依据，最新有效约束必须继续生效。重执行生成新的尝试标识和产物版本，旧结果保留失效标记供审计。对已发送通知、提交订单或写入外部系统的操作，先核对幂等键和外部结果；无法撤销的副作用进入补偿或人工处理路径，不把状态回滚等同于现实世界回滚。

\subsection{面向静默失效与状态丢失的进展检测与统一恢复}

\paragraph{以任务进展补充服务监控}
静默失效检测检查连续窗口内是否出现新的有效产物、已验证事实或待办推进，并结合动作重复、相似输出、资源重复访问和质量返工记录定位停滞。无进展只是触发核验的信号：数据下载、长计算等步骤应使用心跳、进度计数或任务专用检查器，避免把正常等待误判为死循环。模型自报置信度仅作为辅助信号。

检测结果分为继续观察、补充证据、修订本步、局部重规划和暂停。状态中记录触发证据、处理动作及尝试次数；每类恢复设置有限预算。持续返回状态码200而重复生成相同无效产物，仍会因缺少可核验进展进入恢复路径。

\paragraph{统一的状态检查点}
账本、图状态和执行位置在逻辑上组成一个可恢复状态单元：
\begin{equation}
K_t=\langle S_t,\mathcal{E}_t,\mathcal{Q}_t,
\mathcal{R}_t^{\mathrm{mem}},\mathcal{U}_t,\nu_t\rangle,
\end{equation}
其中，$\mathcal{E}_t$为执行前沿，$\mathcal{Q}_t$为消息及接收版本恢复引用，$\mathcal{R}_t^{\mathrm{mem}}$为本次运行使用的记忆版本清单，$\mathcal{U}_t$为外部副作用和幂等记录，$\nu_t$为提交标识。无需复制整个跨项目记忆库，但必须能定位当时使用的版本及其后续撤销信息。

持久化时先写入各对象与引用，再写完成清单；只有引用完整且版本一致的清单才被选为恢复点。该协议把现有上下文与图检查点组织为统一恢复入口，底层仍可使用不同存储。并行写入和跨存储故障需要提交协调，不能仅靠相同文件名视为原子提交。

\begin{figure}[htbp]
\centering
\fbox{\parbox{0.92\textwidth}{\centering
选择完整检查点 $\rightarrow$ 核对状态版本与执行前沿\\[3pt]
$\rightarrow$ 合并最新有效约束、失效记录与外部结果\\[3pt]
$\rightarrow$ 重建消息接收版本、按需唤醒记忆\\[3pt]
$\rightarrow$ 检查依赖、预算与资源 $\rightarrow$ 继续或局部重执行}}
\caption{版本化共享状态的恢复与继续执行流程}
\label{fig:anchor_recovery}
\end{figure}

\subsection{核心算法、实现边界与验证口径}

\begin{algorithm}[htbp]
\caption{共享状态与四级记忆驱动的节点执行}
\label{alg:anchor_runtime}
\KwIn{共享状态$S_t$，节点$n$，四级记忆$M$，恢复与预算策略$b$}
\KwOut{新状态版本、恢复请求或暂停原因}
读取$n$所需的状态版本、任务目标及有效约束\;
根据任务、角色和权限检索记忆，核验来源、有效期与撤销状态\;
构造任务锚点、状态投影、当前输入和相关记忆\;
对最终上下文检查预算；必要信息超限时拆分或暂停\;
执行$n$，保留候选产物、输入依赖和工具记录\;
\eIf{验收通过且依赖版本仍有效}{
    提交状态增量；对已完成工作生成候选记忆并执行准入\;
    保存一致的状态、执行位置与必要恢复引用\;
}{
    隔离候选产物，定位错误依赖并计算受影响范围\;
    在恢复预算内补证、局部重执行或暂停\;
}
检查实际进展并记录停滞与处理动作\;
\end{algorithm}

\begin{algorithm}[htbp]
\caption{带有效性检查的分级证据唤醒}
\label{alg:anchor_wakeup}
\KwIn{查询$q$，状态$S_t$，四级记忆$M$，配置$p$}
\KwOut{证据集合$E$及依赖版本清单}
按权限、作用域、时间和验证状态过滤候选记忆\;
收集词项与语义候选，统一重排、去重，选择至多$K$条\;
对可能冲突或需要核对参数的记录展开原文\;
隔离已撤销、来源失效或适用条件不满足的记录\;
按剩余输入预算组织证据，保留来源与版本\;
\Return $E$及依赖版本清单\;
\end{algorithm}

设候选记忆数为$N$，展开条数为$K$、单条组织长度为$R$，检索成本为$T_{\mathrm{ret}}(N)$，则完整排序路径为
\begin{equation}
T_{\mathrm{wake}}=T_{\mathrm{ret}}(N)+O(N\log N+KR).
\label{eq:anchor_wakeup_cost}
\end{equation}
预算计数随装配文本长度增长；失效传播在受影响依赖子图上遍历，成本为$O(|V_D|+|E_D|)$。检查点写入开销取决于状态和增量载荷大小。以上不包含模型、原文读取及外部工具的实际耗时。

现有工程已提供账本、预算控制、分层记忆、归档、漂移检测及上下文和图检查点等基础组件。本节进一步给出统一版本清单、验收后记忆准入、细粒度依赖撤销及局部重执行的目标设计；这些扩展需要在运行时补齐后单独验证，现有组件测试不能证明整个目标架构已实现。

四类问题应分别通过上下文用量与关键证据保留、约束遵守与目标保持、错误下游采纳率与修复范围、静默停滞检出与一致恢复结果评价。仅减少输入长度不足以证明记忆有效，仅检测到异常也不足以证明已经恢复。验证应包含A到B到C的错误记忆注入、正常长计算、并行写入冲突、中断恢复及外部副作用核对。

